%%%% RESUMO
%%
%% Apresentação concisa dos pontos relevantes de um texto, fornecendo uma visão rápida e clara do conteúdo e das conclusões do
%% trabalho.

\begin{resumoutfpr}%% Ambiente resumoutfpr

A demanda por processadores mais eficientes, econômicos e de alto desempenho apresenta um crescimento contínuo, tornando essenciais as pesquisas que exploram modificações arquiteturais voltadas à eficiência energética ou à capacidade de processamento. Contudo, o design de hardware possui barreiras de entrada elevadas, associada a alta complexidade e custos expressivos. Nesse contexto, este trabalho apresenta o planejamento e a especificação para a implementação de um núcleo processador baseado na ISA RISC-V, utilizando especificamente o conjunto de instruções básicas inteiras de 32 bits (RV32I) em ciclo único com arquitetura de Harvard. Para a validação física do circuito, será adotado uma matriz de portas programáveis (FPGA) AMD Artix-7, integrada em uma placa de desenvolvimento da fabricante Digilent, especificadamente a Nexys A7, visando estabelecer uma plataforma didática e modular que facilite investigações futuras e customizações para aplicações específicas. Para fundamentar a proposta, realizou-se uma revisão bibliográfica sobre o desenvolvimento de processadores, identificando o crescimento do ecossistema RISC-V em frentes que englobam desde a computação de alto desempenho e aceleradores de Inteligência Artificial até microcontroladores de baixo consumo. Adicionalmente, revisaram-se os conceitos estruturais da Arquitetura de Harvard, seus componentes internos essenciais e os modelos de processamento em ciclo único, multiciclo e segmentado (\textit{pipeline}). A metodologia proposta estrutura-se de forma sistemática em sete etapas consecutivas, cobrindo desde o estudo e a modelagem modular dos blocos em linguagem de descrição de hardware SystemVerilog até os testes finais em hardware. Por fim, o particionamento do sistema em cinco blocos funcionais independentes e o cronograma estabelecido servem como base para a organização do projeto. Como resultado deste TCC 1, apresenta-se o planejamento e o mapeamento inicial da arquitetura, o que ajuda a guiar e avaliar a viabilidade técnica para a futura implementação prática e testes em bancada na etapa seguinte do trabalho.

\end{resumoutfpr}

